\documentclass[11pt,a4paper]{report}
\usepackage[margin=1in]{geometry}
\usepackage{hyperref}
\usepackage{enumitem}
\usepackage{graphicx}
\usepackage{xcolor}
\usepackage{booktabs}
\usepackage{amsmath}

% -----------------------------------------------------
% bvraghav-tiet-ucs503p-article.sty
%
% NOTE: our template repo does not ship a `tietreport` class.
% This document reuses the same title-block macros already proven
% to compile in project-proposal/main.tex, on top of the `report`
% class (rather than `article`) so \chapter is available.
% -----------------------------------------------------

% Variable: Lab Instructor
\makeatletter \newcommand{\BvrSetLabInstructor}[1]{%
  \gdef\Bvr@LabInstructor{#1}}
% default
\newcommand{\Bvr@LabInstructor}{Dr. Raghav %
  B. Venkataramaiyer}
% public accessor
\newcommand{\BvrLabInstructorValue}{\Bvr@LabInstructor}
\makeatother

% Add subtitle
\usepackage{titling}
% -- Set title bold
\pretitle{\begin{center}\bfseries\fontsize{18bp}{18bp}\selectfont}
\makeatletter
\newcommand{\BvrSubtitle}[1]{%
  \posttitle{%
    \par\end{center}
    \begin{center}\large#1\end{center}
    \vskip 0.5em}%
}
\makeatother

% Hints in the section title(s)
\newcommand{\BvrHint}[1]{%
  {\normalfont\normalsize\itshape\hfill #1}}

% -----------------------------------------------------

\usepackage[
  style=alphabetic,
  backend=biber,
  sorting=nty,
  maxnames=5,
  minnames=3,
  doi=false,
  isbn=false,
  url=false,
]{biblatex}
\addbibresource{references.bib}

% TODO: confirm all three email addresses before submitting.

\title{Hostel Room Swapping Service}

\BvrSetLabInstructor{Dr. Raghav B. Venkataramaiyer}

\BvrSubtitle{%
  A Multi-Party Hostel Room Exchange Matching Service \\[0.3em]
  Project Report --- Prototype Stage, UCS503P \\
  Group: \texttt{3C51} \\
  Submitted to: \BvrLabInstructorValue \\ \bfseries
  Thapar Institute of Engineering and Technology%
}

\author{
  Aabhas Khandelwal, CSED%
  \thanks{Roll No: \texttt{1024030673}, %
    Email: \texttt{akhandelwal2\_be24\@thapar.edu}}
  \and
  Aakash Chawla, CSED%
  \thanks{Roll No: \texttt{1024031097}, %
    Email: \texttt{achawla1\_be24\@thapar.edu}}
  \and
  Joy Soni, CSED%
  \thanks{Roll No: \texttt{1024030674}, %
    Email: \texttt{jsoni\_be24\@thapar.edu}}
}

\date{\today}

\begin{document}
\maketitle
\tableofcontents

% ==================================================
% CHAPTER 1: INTRODUCTION
% ==================================================

\chapter{Introduction}

\section{Background and Context}

Hostel room allotment assigns each student one room for the semester,
but that allotment is fixed once made. A student who ends up on the
wrong floor, facing the wrong direction, or away from friends has no
structured way to change it. The informal remedy --- asking around,
posting in a group chat --- only ever surfaces a handful of offers,
and a workable trade almost always needs more than two people.

A direct one-to-one swap requires two students to each want exactly
what the other already has. That coincidence is rare. What usually
does exist is a longer exchange: student A takes B's room, B takes
C's, and C takes A's, leaving everyone better off. No single student
can see that chain from inside their own conversation, because seeing
it requires knowing every other student's current room and every
other student's preference at once.

This project treats room allocation as a visibility problem rather
than a negotiation problem. The service holds every student's current
room and every student's stated preference centrally, computes the
exchange cycles hidden inside that data, and presents them back as
concrete, executable swaps.

\section{Problem Statement}

Framed as requirements, a hostel room exchange service should:

\begin{enumerate}
    \item Accept a student's preference for a specific room, with a
    defined fallback --- floor, then direction --- when that exact
    room cannot be obtained.

    \item Compute, from every participant's current room and stated
    preference simultaneously, the assignment of maximum total
    satisfaction across the group, not the loudest individual request.

    \item Generate genuinely different alternative assignments, not
    one answer marginally degraded, so that a shorter and more easily
    approved exchange can be offered alongside the mathematically
    optimal one.

    \item Guarantee that no participant in an offered exchange ends up
    worse off than their starting room, as a structural property of
    the system rather than an incidental outcome.

    \item Require every member of an exchange chain to individually
    approve it before any room actually changes, and ensure one
    member's refusal cannot silently affect any other exchange running
    in the same round.

    \item Preserve a complete, non-rewritable history of who held
    which room and when, so that any executed swap is auditable after
    the fact.
\end{enumerate}

\section{Project Objectives}

\begin{enumerate}

    \item \textbf{Closed-pool assignment.} Model the swap problem as a
    permutation of students over the rooms they already collectively
    hold, so that swap chains are not searched for separately but fall
    directly out of the assignment's cycle structure.

    \item \textbf{Optimal and alternative assignment.} Solve for the
    single best assignment using a proven combinatorial algorithm,
    and generate a bounded number of alternative assignments so that
    a shorter, more executable chain can be preferred when it exists
    at no cost to overall match quality.

    \item \textbf{Feasibility filtering.} Discard any candidate
    exchange that would leave a participant worse off, independent of
    whatever the raw optimisation objective would otherwise select.

    \item \textbf{A real approval workflow.} Model proposal,
    approval, execution, rejection and expiry as an explicit state
    machine, separate from the state machine that governs a matching
    round itself.

    \item \textbf{Persistent, auditable storage.} Store allocation
    history as an append-only log in a managed relational database,
    with database-level enforcement of that guarantee.

    \item \textbf{Measurable evaluation.} Evaluate the engine against
    a reproducible cohort using metrics that describe the whole
    participant group, not only the students who moved.

\end{enumerate}

% ==================================================
% CHAPTER 2: METHODOLOGY AND DESIGN
% ==================================================

\chapter{Methodology and Design}

\section{System Architecture}

The service is organised as a pipeline connecting a browser-based
interface, a stateless API layer, a pure matching engine, and a
managed relational database.

The major components are:

\begin{enumerate}
    \item Admin and student web pages
    \item FastAPI backend and session-based authentication
    \item Preference and bed-slot modelling
    \item Cost-matrix construction and scoring
    \item Hungarian-algorithm assignment
    \item Murty's-algorithm alternative generation
    \item Cycle decomposition and Pareto feasibility filtering
    \item Approval (proposal) state machine
    \item PostgreSQL persistence with append-only allocation history
\end{enumerate}

The registration pipeline is:

\begin{center}
\textbf{
Student Login
$\rightarrow$
Round Enrolment
$\rightarrow$
Weighted Priorities
$\rightarrow$
Stored Preference Set
}
\end{center}

The matching pipeline is:

\begin{center}
\textbf{
Round Locked
$\rightarrow$
Cost Matrix
$\rightarrow$
Hungarian Assignment
$\rightarrow$
Murty $K$-Best
$\rightarrow$
Cycle Decomposition
$\rightarrow$
Pareto Filter
$\rightarrow$
Proposals Offered
}
\end{center}

Separating registration from matching, and matching from approval,
means each stage can be developed, tested, and reasoned about
independently: the round records what the engine found to be
\emph{possible}, and the proposal state machine records what the
participants actually \emph{agreed to}.

\subsection{Web Interface}

Two role-scoped web views sit in front of the service: a student view
for enrolling in a round, setting priorities, and responding to
proposed exchanges, and a warden (admin) view for creating and locking
a round for a given hostel. Both are served as plain HTML/JavaScript
pages by the same backend that exposes the API, so no separate
frontend deployment is required at this stage. A migration to a
React-based frontend is planned; because the pages consume the same
JSON API a script or a framework component would, that migration does
not require any change to the backend or matching engine.

\subsection{FastAPI Backend}

The backend exposes 27 endpoints across four groups: authentication
(login, logout, session lookup), student actions (enrol, set
priorities, view results, offer or respond to a proposal), admin
actions (create, open, and lock a round; declare room-type inventory),
and two diagnostic endpoints (\texttt{/api/schema}, which reports the
live database structure, and \texttt{/api/demo}, a fixed-cohort
endpoint used for reproducible evaluation, described in
Section~\ref{sec:eval-methodology}). Every write endpoint that
requires a specific role is rejected at the backend if the
authenticated session does not hold that role.

\subsection{Bed-Slot Modelling}
\label{sec:bed-slot}

A student's stated preferences --- hostel, floor, direction, room
type, washroom --- are properties of a \emph{room}. Assignment,
however, has to happen over a \emph{one-to-one} structure for the
Hungarian algorithm to apply at all. A shared room with three beds is
not one assignable unit; it is three. The system therefore assigns
students to individual bed slots rather than rooms directly, with each
room expanding into as many slots as its declared capacity. This
keeps the assignment problem a clean permutation --- as many slots as
participating students --- while still allowing shared rooms and
mutual roommate requests to be expressed, without which the whole
Hungarian/Murty formulation would not apply.

\section{Preference Scoring}

Each student's satisfaction with a candidate room is the fraction of
their stated preference weight that room satisfies. Preferences are
distinguished as hard or soft: an unmet soft preference lowers the
score, while an unmet hard preference removes the room from
consideration entirely rather than merely penalising it. In the
current evaluation cohort, priorities follow a fixed fallback ladder
--- a named room is worth more than its floor, which is worth more
than its direction --- so that a student who cannot obtain their exact
room still scores well for landing near where their friends are.

\section{Matching Algorithm}

The pool of participants in a round is closed: the only rooms
available are the ones the participants already collectively hold.
Consequently, any solution is a permutation of students over slots,
and every permutation decomposes uniquely into disjoint cycles. Swap
chains are therefore not searched for separately from the assignment
--- they are already present inside it, and the system reads them out.

\subsection{Assignment Optimisation}

The Hungarian algorithm, via
\texttt{scipy.optimize.linear\_sum\_assignment}
\autocite{kuhn1955hungarian,virtanen2020scipy}, solves the
students-by-slots cost matrix for the assignment of maximum total
satisfaction in $O(n^3)$.

\subsection{Alternative Generation}

Murty's algorithm \autocite{murty1968ranking} enumerates the $K$ best
assignments by repeatedly forbidding one edge of an already-found
assignment and re-solving. This is a correctness requirement rather
than a convenience feature: the single mathematically optimal
assignment frequently demotes one participant to benefit several
others, which the feasibility filter (Section~\ref{sec:feasibility})
then rejects outright --- the alternatives are where an
\emph{executable} exchange is actually found.

\subsection{Swap Feasibility}
\label{sec:feasibility}

Each candidate assignment is decomposed into its disjoint exchange
cycles. Because the cycles are disjoint, they execute independently:
a single student withdrawing from one cycle provably cannot affect
any other cycle in the same assignment. A Pareto filter then discards
any cycle in which a participant would end up worse off than their
current room, keeping only chains that every member has a rational
reason to accept.

\subsection{Alternatives Considered}

Gale--Shapley stable matching \autocite{gale1962college} was
considered and rejected: it assumes two-sided preferences, where both
sides rank each other, and rooms hold no preference over students, so
half the algorithm has no input. Top Trading Cycles
\autocite{shapley1974cores} fits the exchange shape of the problem
directly, but operates over a strict ranked list of items and cannot
represent preferences that are weighted across several attributes at
once, as ours are.

\section{Approval Workflow}

A candidate exchange chain is issued as a proposal, not applied
immediately. Every member of the chain must individually approve it;
a single rejection kills only that one proposal. A round moves through
its own lifecycle --- draft, open, locked, running, completed ---
independently of the approval state of the proposals it produced,
which is what allows a round to be complete (the engine has run and
produced its answer) while individual proposals inside it are still
being negotiated.

\section{Persistence}

The schema spans 14 tables in a managed PostgreSQL database
(Supabase): students and sessions; hostels and their declared
room-type inventory; rooms and bed slots; the allocation history;
preference sets; swap rounds and enrolments; and the proposal and
swap-chain-member tables backing the approval workflow. Allocation
history is append-only, enforced by a database trigger rather than by
application-level convention, so a defect in a future client cannot
silently rewrite who held which room and when.

\section{Technical Stack}

\begin{table}[H]
\centering
\begin{tabular}{p{0.32\textwidth} p{0.53\textwidth}}
\toprule
\textbf{Component} & \textbf{Technology} \\
\midrule
Programming language & Python \\
Backend framework & FastAPI \\
Assignment solver & SciPy (\texttt{linear\_sum\_assignment}) \\
Alternative generation & Murty's algorithm (in-house, over SciPy) \\
Database & PostgreSQL (Supabase, managed) \\
Authentication & Session cookies, server-side session table \\
Testing & pytest \\
CI/CD & GitHub Actions \\
Frontend (current) & Server-rendered HTML / JavaScript \\
Frontend (planned) & React \\
\bottomrule
\end{tabular}
\caption{Technical Stack}
\label{tab:techstack}
\end{table}

% ==================================================
% CHAPTER 3: IMPLEMENTATION AND EVALUATION
% ==================================================

\chapter{Implementation and Evaluation}

\section{Implementation}

The matching engine is implemented as a pure Python package with no
framework or database import anywhere in it, so it can be exercised
and tested with no infrastructure running at all.

\begin{itemize}
    \item \texttt{scoring.py} --- preference weighting and match
    calculation.
    \item \texttt{matrix.py} --- students-by-slots cost matrix
    construction.
    \item \texttt{solver.py} --- Hungarian-algorithm assignment.
    \item \texttt{kbest.py} --- Murty's $K$-best alternative search.
    \item \texttt{cycles.py} --- cycle decomposition into exchange
    chains.
    \item \texttt{feasibility.py} --- the Pareto feasibility filter.
    \item \texttt{ranking.py} --- candidate scoring and selection.
    \item \texttt{proposal.py} --- the approval state machine.
    \item \texttt{rounds.py} --- the round lifecycle and orchestration
    of the pipeline above.
    \item \texttt{db.py} --- the only module that speaks SQL; loads
    and persists the same data shapes the engine consumes.
    \item \texttt{auth.py} / \texttt{api.py} --- session
    authentication and the FastAPI route layer.
\end{itemize}

\subsection{Round Lifecycle}

A round for one hostel moves through \texttt{draft}
$\rightarrow$ \texttt{open} $\rightarrow$ \texttt{locked}
$\rightarrow$ \texttt{running} $\rightarrow$ \texttt{completed}.
Locking a round is a one-way transition: it closes registration and
triggers the matching pipeline in the same action, after which the
round cannot return to \texttt{open}.

\subsection{Proposal Lifecycle}

Each offered exchange chain moves through \texttt{proposed}
$\rightarrow$ \texttt{approved} $\rightarrow$ \texttt{executed}, or
exits via \texttt{rejected} or \texttt{expired}. Execution appends new
allocation rows for every member of the chain; no existing row is ever
modified.

\section{Evaluation Methodology}
\label{sec:eval-methodology}

No public benchmark exists for hostel room-swap matching, so the
engine is evaluated against a reproducible synthetic cohort rather
than an external dataset. A pool of 30 students, one per room, in a
single hostel, was drawn once from a seeded random generator and
frozen, so evaluation results are exactly reproducible run to run
rather than depending on which random cohort happened to be sampled
that day. Each student's priorities follow the room/floor/direction
fallback ladder described in Section~\ref{sec:bed-slot}, with no
preference toward the hostel itself, since every student in the
cohort already lives there.

The cohort is exercised through the service's own \texttt{/api/demo}
endpoint, which runs the identical matching pipeline used by a real
locked round, so the benchmark evaluates the engine's real code path
rather than a simulation of it.

\section{Evaluation Metrics}

\subsection{Cohort Preference Satisfaction}

The primary metric is the mean weighted match score across
\emph{every} participant in the cohort, computed before and after one
matching round --- not only across the students who moved. This is a
deliberate choice: averaging over movers alone would rank a two-person
swap in which both reach 100\% ahead of an eleven-person rotation that
lifts the entire cohort, which is not the outcome the service exists
to produce.

\subsection{Exact-Room, Floor, and Direction Fulfilment}

Secondary metrics report the proportion of the cohort receiving the
exact room they named, the proportion receiving at least their
requested floor, and the proportion receiving at least their requested
direction --- the three rungs of the fallback ladder.

\subsection{Participants Left Worse Off}

This count must be zero by construction, since the Pareto filter
(Section~\ref{sec:feasibility}) discards any chain violating it; it is
reported as a measured value rather than assumed, to confirm the
filter behaves as specified on real data rather than only in unit
tests.

\subsection{Longest Simultaneous-Approval Chain}

The size of the largest exchange cycle in the offered assignment,
reported because it is the practical execution risk of a chain: every
additional member is one more independent approval that must be
obtained before the whole chain can execute.

\section{Experimental Configurations}

Two configurations of the same pipeline were evaluated on the frozen
cohort, differing only in $K$, the number of alternative assignments
Murty's algorithm explores before the best candidate is selected:

\begin{enumerate}
    \item \textbf{$K = 1$:} only the single mathematically optimal
    assignment is considered.
    \item \textbf{$K = 80$ (production):} the full Murty's search is
    run, and the assignment maximising cohort satisfaction is
    selected from among the alternatives, breaking ties toward a
    shorter longest chain.
\end{enumerate}

\section{Results}

\begin{table}[H]
\centering
\begin{tabular}{lcc}
\toprule
\textbf{Metric} & \textbf{$K=1$} & \textbf{$K=80$} \\
\midrule
Cohort satisfaction, before $\to$ after & 11\% $\to$ 77\% & 11\% $\to$ 77\% \\
Students moved & 28 / 30 & 28 / 30 \\
Received exact requested room & 20 / 30 & 20 / 30 \\
Participants left worse off & 0 & 0 \\
Longest simultaneous-approval chain & 28 & 11 \\
\bottomrule
\end{tabular}
\caption{Matching Engine Evaluation, Frozen 30-Student Cohort}
\label{tab:evaluation}
\end{table}

\section{Analysis of Results}

\subsection{Match Quality Is Unaffected by $K$}

Cohort satisfaction, the count of students moved, and the count
receiving their exact requested room are identical at $K=1$ and
$K=80$. On this cohort, exploring alternative assignments does not
trade away any match quality at all.

\subsection{Executability Is Not}

The longest chain requiring simultaneous approval falls from 28
students to 11 between the two configurations, at no cost measured
above. A chain of 28 means 28 independent people must each agree
before anything executes, and a single hold-out anywhere in that chain
blocks the entire exchange. A chain of 11 carries proportionally less
of that risk, and because the underlying cycles are disjoint by
construction (Section~\ref{sec:feasibility}), a shorter maximum chain
generally also means the assignment is expressible as several smaller
independent chains rather than one large one.

\subsection{Zero Worse Off, Confirmed on Real Data}

The Pareto guarantee is a property proved of the filter's logic in
unit tests; Table~\ref{tab:evaluation} confirms it also holds when the
filter is exercised against a full 30-student assignment rather than a
small hand-constructed case.

\section{Comparative Analysis}

\begin{table}[H]
\centering
\begin{tabular}{p{0.30\textwidth} p{0.55\textwidth}}
\toprule
\textbf{Observation} & \textbf{Interpretation} \\
\midrule

Satisfaction identical at $K=1$ and $K=80$ &
Murty's alternative search finds a more executable assignment without
sacrificing overall match quality on this cohort. \\

Longest chain: 28 at $K=1$, 11 at $K=80$ &
Far fewer simultaneous approvals are required for the same outcome,
directly reducing the practical risk of the exchange collapsing before
it executes. \\

20 of 30 receive their exact requested room &
The majority of demand is satisfiable without falling back to floor
or direction on this cohort. \\

0 participants worse off, in both configurations &
The feasibility filter's guarantee held with no exception across the
full assignment, not only on average. \\

\bottomrule
\end{tabular}
\caption{Comparison of Evaluated Configurations}
\label{tab:comparison}
\end{table}

These results are the concrete justification for including Murty's
algorithm in the design at all: a single Hungarian solve would have
been sufficient if match quality were the only concern, but it is not
--- an assignment that cannot realistically be approved is not a
useful answer, and $K$-best search is what turns an optimal-on-paper
assignment into one that is also practical to execute.

\section{Software Engineering Process}

The full automated test suite --- 127 tests --- runs in approximately
two seconds with no live database, a direct consequence of the
matching engine importing no framework or database driver: unit and
scenario tests exercise the pure Python pipeline directly, and only a
small, isolated set of tests exercises the PostgreSQL integration path
so the rest of the suite remains hermetic on a machine with no
database credentials at all. GitHub Actions runs the suite on every
push and pull request, across Python 3.10, 3.11, and 3.12.

Three defects surfaced only once real infrastructure was in the loop,
each caught before it reached a shared environment:

\begin{itemize}[leftmargin=0.7cm]
    \item A preference criterion added to the Python model was never
    added to the corresponding PostgreSQL enum, which would have
    rejected the very first insert against a fresh database.
    \item A database driver installed by hand on a development
    machine was never declared as a project dependency, so it worked
    locally and failed immediately in a clean CI environment.
    \item A deliberate offline fallback --- used so a demo cannot be
    broken by a database outage --- initially made a genuine
    misconfiguration indistinguishable from a correctly working
    system, until the API response was made to report which data
    source it had actually used.
\end{itemize}

% ==================================================
% CHAPTER 4: CONCLUSION
% ==================================================

\chapter{Conclusion}

\section{Summary of Work}

The prototype implements a complete pipeline from student preference
submission through to an executed room swap: a pure matching engine
combining the Hungarian algorithm with Murty's $K$-best alternative
search and a Pareto feasibility filter; a proposal-based approval
workflow layered on top of a separate round lifecycle; and a
PostgreSQL-backed persistence layer with an enforced append-only
allocation history. A session-authenticated web interface exposes
both a student role and a warden (admin) role over the same backend.

\section{Key Findings}

\begin{enumerate}

    \item \textbf{The engine finds exchanges no participant could see
    alone.} Multi-party chains are read directly out of a computed
    permutation rather than searched for separately.

    \item \textbf{Alternative-assignment search reduces execution
    risk at no measured cost to match quality.} On the evaluated
    cohort, raising $K$ cut the longest simultaneous-approval chain
    from 28 students to 11 while leaving cohort satisfaction
    unchanged.

    \item \textbf{The feasibility guarantee held with zero
    exceptions} across the full evaluated assignment, not only in
    isolated unit tests.

    \item \textbf{The system performs practically.} The evaluated
    30-student assignment computes in well under a second, and the
    matching engine's freedom from framework and database dependencies
    keeps the automated test suite fast enough to run on every push.

    \item \textbf{Testing against real infrastructure caught real
    defects.} Three separate issues --- a schema/model mismatch, an
    undeclared dependency, and a silent fallback --- were each found
    because the prototype was exercised against a genuine database
    and CI pipeline early, rather than only against hand-written unit
    tests.

\end{enumerate}

\section{Limitations}

\begin{itemize}[leftmargin=0.7cm]

    \item The engine has been evaluated against a seeded synthetic
    cohort, not yet against real student registrations for a live
    hostel.

    \item Shared (multi-bed) rooms and mutual roommate requests are
    supported by the bed-slot model and scoring function, but are not
    exercised by the current evaluation cohort, which uses one student
    per room so every chain reads as an unambiguous rotation.

    \item Proposal expiry and round cancellation exist in the state
    model but have not yet been demonstrated end-to-end on a live
    clock.

    \item The frontend is presently server-rendered HTML/JavaScript
    rather than the React frontend named in the project proposal.

\end{itemize}

\section{Future Work}

\begin{enumerate}

    \item Onboard real student preference registrations for at least
    one hostel, replacing the seeded evaluation cohort.

    \item Exercise shared rooms and mutual roommate matching
    end-to-end, including on the live evaluation cohort.

    \item Demonstrate the reject, expiry, and round-cancellation
    paths live rather than only in the test suite.

    \item Migrate the frontend to React against the existing,
    unchanged API contract.

    \item Evaluate the engine at larger cohort sizes to characterise
    how satisfaction, chain length, and solve time scale together.

\end{enumerate}

\section{Final Remarks}

This prototype demonstrates that a hostel's worth of room-swap
requests can be resolved as a single closed assignment problem, with
the exchange chains falling directly out of that assignment's cycle
structure rather than being searched for separately. The evaluation
shows that generating alternative assignments is not merely an
academic completeness exercise: on the tested cohort it produced a
markedly more executable outcome --- an 11-person longest chain in
place of a 28-person one --- without giving up any measured match
quality. The remaining work is primarily about exercising paths the
architecture already supports against real data, rather than
redesigning the underlying approach.

% ==================================================
% BIBLIOGRAPHY
% ==================================================

\printbibliography[heading=bibintoc]

\end{document}
